% LAB BOREDOM — module L3, The Findings. GUIDED DRAFTING (desktop recipe), Section 2 run.
% Session 2026-08-07. Spec: outline v8 §L3 — v7's L3 is VOID (it was organized around film
% separation from L3.1 onward). This module is built on the criterion analysis instead.
%
% ★ GOVERNING RULING (author, 2026-08-07), restated because it determines this module's shape:
%   film-separation statistics are EXCLUDED ENTIRELY — no 99 contrasts, no Friedman/Wilcoxon
%   across films, no Holm-surviving contrast list, no "this measure distinguishes the boring
%   film from the clinical one". The criterion analysis IS the results. Gaze deviation
%   (per-participant baseline) is the principal finding. The Spearman/Pearson comparison chart
%   goes in the BODY as an image; the fuller table for non-load-bearing measures to the appendix.
%   All three films treated equally; the clinical stimulus is not foregrounded.
%
% All figures below read from boredom-analysis-v5 out/criterion-validity.csv AFTER the D-39 and
% D-40 re-runs (2026-08-07). Pooled, n=36 episodes, 12 participants, within-subject z, 20,000
% within-subject permutations.
%   gaze deviation (per-participant)  boredom +0.590 (0.0024) | r +0.623 (0.0016)
%                                     engagement -0.585 (0.0022) | r -0.609 (0.0021)
%   gaze deviation (per-file)         boredom +0.524 (0.0074) | r +0.574 (0.0036)
%                                     engagement -0.527 (0.0075) | r -0.590 (0.0030)
%   eye closure                       boredom +0.475 (0.0173) | r +0.580 (0.0038)
%                                     engagement -0.435 (0.0279) | r -0.535 (0.0073)
%   gaze deviation (device-forward)   boredom +0.359 (0.0751) | engagement -0.391 (0.0515) — neither holds
%   pupil dilation                    boredom -0.172 (0.4037) | engagement +0.036 (0.8608)
%   cognitive load                    boredom -0.092 (0.6516) | engagement +0.107 (0.5954)
%   HR                                boredom +0.088 (0.6774) | engagement +0.036 (0.8647)
% Leave-one-out leverage (2026-08-07): every influential episode SUPPRESSES the Pearson
%   coefficient. Dropping S06/INT raises gaze r from +0.623 to +0.707; dropping S06/BOR raises
%   eye closure r from +0.577 to +0.655. No point props either coefficient up.
%
% BATCH 1 = the principal finding (gaze deviation). Awaiting author review.
% ★ FIGURE PENDING: the Spearman/Pearson comparison chart is called for in this batch.

Christine King and I built this experiment to determine whether boredom, engagement, and
attention span could be identified and tracked in physiological data, using participants' own
reports of those states as the reference each physiological measure would be judged against.
The question has an answer, and the answer is yes: eye tracking accomplishes it. Of the measures considered across the
thirty-six episodes---EEG, pupil diameter, how much of an episode the eyes spent closed, and
where the eye was pointing---the measures taken by the eye tracker are those that move with what
participants reported, and the strongest of them is gaze deviation.

How far a participant's gaze wandered from their own habitual direction rises with the boredom
they reported, at a rank correlation of $+$0.590, and falls with the engagement they reported,
at $-$0.585.\footnote{A correlation summarizes how closely two quantities move together in a
single number between $-$1 and $+$1. A value of 0 means they are unrelated: knowing one tells
nothing about the other. A value of $+$1 means they move together perfectly, so that every
increase in one is matched by a proportionate increase in the other, and $-$1 means they move
perfectly in opposite directions. The sign gives the direction of the relationship and the size
gives its strength, so $+$0.590 and $-$0.585 describe relationships of almost identical
strength running in opposite directions---gaze wandering further as reported boredom rises, and
less far as reported engagement rises.} Neither figure could plausibly have arisen by chance: across twenty thousand
rearrangements in which each participant's ratings were shuffled among their own episodes, an
association this strong appeared in fewer than one run in four hundred for boredom and fewer
than one in four hundred and fifty for engagement. The relationship is also close to a straight
line rather than merely a consistent ordering. Measured as a linear correlation the same
association reaches $+$0.623 and $-$0.609, slightly stronger than the rank figures rather than
weaker, which is what happens when the underlying relationship really is proportional and the
ranking is discarding information rather than protecting against a distortion.

% ****** FIGURE PENDING: Spearman/Pearson comparison chart, body placement.
\begin{center}
\fbox{\parbox{0.9\linewidth}{\centering\textbf{****** FIGURE PENDING ******}\\
The Spearman and Pearson coefficients for every measure, side by side, against both reported
boredom and reported engagement. Generated from
\texttt{out/criterion-validity.csv}; the figure pass requires matplotlib, which was absent from
the environment in which the analysis was last run.}}
\end{center}

That the two coefficients agree matters more than either does alone. A rank correlation
establishes only that a participant's episodes fall in the same order on both measures---that
the film they found dullest is also the film their eye wandered furthest during. A linear
correlation of the same size establishes something stronger: that the wandering scales with the
boredom, so that a participant reporting two additional points of boredom shows roughly twice
the additional deviation of one reporting a single point. Across the whole family of measures
the two coefficients differ by a median of 0.032, and where they differ at all the linear
figure is generally the larger. One further check answers the objection that would otherwise be
obvious---that one or two unusual participants are producing the whole effect. Each participant
was removed in turn, the remaining eleven were standardized again from the beginning, and the
correlation recomputed on them alone.\footnote{Removing a whole participant rather than a single
episode is what the design requires. Each participant's three values are scored against their own
average, so the three are bound together---remove one episode and the two left behind are still
centered on a figure that included it. Removing all three of someone's episodes takes out a
self-contained unit and leaves nothing behind that depends on them.} Had the relationship
depended on a few people, dropping one of them would have collapsed it. It never does. Across
all twelve recomputations the correlation stays between $+$0.589 and $+$0.755 against reported
boredom, and between $-$0.574 and $-$0.728 against reported engagement. The two largest movements
are upward---removing either of two participants raises the correlation rather than lowering
it---which means those participants were holding the relationship down rather than creating it.
The association does not rest on anyone in particular.

Two properties of this measure are worth stating before anything is built on it. It is measured
against each participant's own habitual gaze direction across their three films, which is what
allows it to register that someone held their eye differently during one film than another
without also registering how the headset happened to sit on their face. It is also silent about
what any particular film does: a participant who found the tutorial absorbing and the clinical
footage tedious contributes to this association exactly as strongly as one who found the
reverse, because the measure is scored against what that person reported and not against which
film they were watching.

% ============================================================================
% BATCH 2 = eye closure as the second measure, then the sleep-fight results.
% Drafted 2026-08-07. Verified figures (out/criterion-validity.csv, post D-39/D-40):
%   eye closure vs boredom    rho +0.475 (perm 0.0173) | r +0.580 (perm 0.0038)
%   eye closure vs engagement rho -0.435 (perm 0.0279) | r -0.535 (perm 0.0073)
%   pupil dilation vs boredom -0.172 (0.4037) | vs engagement +0.036 (0.8608)
% Sleep-fight (computed 2026-08-07, NOT in outline v7; recorded in outline v8):
%   vs boredom    rho +0.771 | r +0.823, within-subject permutation p < 0.0001
%   vs engagement rho -0.752 | r -0.782, within-subject permutation p < 0.0001
%   leave-one-out: every influential episode SUPPRESSES r (drop S06/BOR -> +0.873)
%   ★ WITHIN-PARTICIPANT PAIRED CONTRAST: each participant's OWN most-boring episode vs their
%     OWN most-engaging episode. 11 of 12 fought sleep harder on their most-boring episode,
%     0 the other way, 1 tied (S06, 6/6). Exact sign test two-sided p = 0.0010 = THE FLOOR at
%     11 pairs. Exact sign-flip permutation agrees at 0.0010.
%   ★ MANDATORY SHARED-METHOD CAVEAT (outline v8): sleep-fight is a survey item correlated with
%     other survey items -- same person, same instrument, same moment. It evidences the survey's
%     internal coherence, NOT independent corroboration from outside the survey. State it every
%     time the figure appears.
%   S04 detail: their most-boring episode was CLC (bored 9), most-engaging was BOR (engaged 7);
%     sleep-fight 6 against 1. Follows the pattern by one of the larger margins in the set.
% ============================================================================

The second measure to move with what participants reported is how much of an episode they spent
with their eyes closed. It rises with reported boredom at a rank correlation of $+$0.475 and
falls with reported engagement at $-$0.435, and as a linear correlation it reaches $+$0.580 and
$-$0.535. Both hold up against the same shuffling test, appearing by chance in fewer than two
runs in a hundred for boredom and fewer than three in a hundred for engagement. The measure is
cruder than gaze deviation and it behaves accordingly: closing the eyes is an all-or-nothing act
that happens intermittently, where gaze wanders continuously, so a single measure of closure
compresses a whole episode into one proportion. That it nonetheless tracks what participants
said, and does so on both questions, is the point. Two independent things the eye does---where
it points and whether it is open---move together with the same reports.

A third measure of the same experience comes from the survey itself. Participants were asked
after each film how hard they had fought to stay awake, and the answers track their reported
boredom closely---a rank correlation of $+$0.771, and $-$0.752 against engagement, both so far
beyond chance that no rearrangement in twenty thousand came near them. That is a stronger
association than either eye measure produces, and the reason it is stronger is also the reason
it counts for less. The sleep question and the boredom question were answered by the same person,
about the same film, in the same breath, on the same sheet. An association between two items on
one questionnaire shows that the questionnaire is internally coherent. It cannot show that
anything outside the questionnaire agrees with it, which is what the eye measures show and what
this measure, by construction, cannot.

What the sleep item can do is answer an objection to itself, and the answer is the more useful
result. Drowsiness is not only a response to a film. A participant who slept badly the night
before will fight sleep through all three episodes, and that has nothing to do with boredom. The
way to remove it is to let each participant serve as their own control: take the episode that
participant rated most boring, take the episode the same participant rated most engaging, and
compare the drowsiness they reported in each. Whatever tiredness they carried into the room is
present on both sides and cancels. Eleven of the twelve participants fought sleep harder during
their own most-boring episode than during their own most-engaging one. No participant reported
the reverse---not one person in the study fought sleep harder during the episode they had called
their most engaging. The twelfth gave the same figure for both. The probability of that can be worked out without
any apparatus. If drowsiness had nothing to do with which episode a participant found boring,
each of the eleven would be a coin toss, and eleven tosses landing the same way happens once in
2,048 times; doubling that to allow the run to have gone either direction gives about one chance
in a thousand. With twelve participants that is also the strongest result the sample size can
express, since a perfect split is the most lopsided outcome available.

% ---------------------------------------------------------------------------
% FATIGUE-AFTER INTEGRATION (computed 2026-08-07; author: "great finding, integrate it with
% the rest of the results and clearly link it to the other data sets").
% The seventh survey item WAS collected and IS in the tree, but had only ever been tested for
% film separation. Scored against what participants reported for the first time here:
%   fatigue AFTER  vs boredom    rho +0.640 (perm 0.0010) | r +0.709 (0.0004)
%                  vs engagement rho -0.616 (perm 0.0019) | r -0.670 (0.0007)
%   depletion      vs boredom    rho +0.516 (0.0088)      | r +0.549 (0.0055)
%                  vs engagement rho -0.420 (0.0355)      | r -0.458 (0.0234)
%   ★ fatigue BEFORE vs boredom  rho +0.071 (0.7229)      -- THE CONTROL, a clean null
%                  vs engagement rho -0.208 (0.3031)
%   within-participant (own most-boring vs own most-engaging): fatigue after 10 higher /
%     1 lower / 1 tied, exact sign test p = 0.0117. Depletion 9/2/1, p = 0.0654 (n.s.).
% ★ THE INTEGRATIVE ARGUMENT this passage exists to make: the fatigue-BEFORE null is what
%   defends the survey items against the shared-method objection. If participants were simply
%   telling a consistent story across one questionnaire, the before-item would correlate too.
%   It does not. The survey discriminates between what a participant brought into the room and
%   what the film did to them, which is a discrimination a mere response-style artifact cannot
%   produce. Link forward to the eye measures, which supply the same conclusion from outside
%   the survey entirely.
% ---------------------------------------------------------------------------

The survey's last item completes the picture. Participants were asked how tired they felt after each film as well as
before it. Fatigue reported afterward tracks reported boredom at $+$0.640 and reported
engagement at $-$0.616, which makes it the strongest single correlate of boredom anywhere in
this study---ahead of gaze deviation, ahead of eye closure, ahead of the sleep item. The
within-participant test agrees: ten of the twelve were more tired after their own most-boring
episode than after their own most-engaging one, one reported the reverse, and one gave the same
figure for both.\footnote{The test is an exact sign test: of the eleven participants who did not give
identical figures for both episodes, count how many went each way, then calculate directly how
often a split that lopsided would occur if the direction were decided by chance. It occurs
about once in eighty-five times, a two-sided probability of 0.0117. The same comparison run on
depletion---the difference between the fatigue reported before a film and after it---points the
same way but less sharply, at $+$0.516 against reported boredom, with nine of twelve in the
expected direction and a probability of 0.065, which falls short of the 0.05 conventionally
treated as the threshold for a result. That is what should be expected. Depletion is built by
subtracting one reported number from another, and subtracting two imprecise figures yields a
figure less precise than either, so the same underlying relationship shows through it more
faintly.}

The item asked before each film is what turns this from a fourth correlation into an argument.
Fatigue reported \textit{before} watching has no relationship to how boring that film was
subsequently rated---$+$0.071 against boredom, and $-$0.208 against engagement, neither
remotely distinguishable from chance. Someone constructing a consistent account of a tedious
afternoon has no reason to exempt the question about how they felt walking in. It is exempted.
The survey distinguishes what a participant carried into the room from what the film did to
them.

Read together, the four measures converge from two directions that do not share a source. From
inside the survey, three separate questions---how bored, how hard the fight against sleep, how
tired afterward---move together, while a fourth question about the state preceding the film
stays flat. From outside it, two measurements taken by an eye tracker that has no access to any
of those answers move with them: gaze wandering further from its own center as reported boredom
rises, and eyes spending more of the episode closed. The survey establishes that participants
are reporting something coherent and specific to the episode rather than a
standing frame of mind they arrived with. The eye
measures establish that whatever they are reporting also has a bodily expression that can be
recorded without asking. Neither alone would carry the claim. Together they describe boredom as
something that shows itself in what a person says, in how long they can hold their gaze, and in
how tired they are when it ends.

% ---------------------------------------------------------------------------
% CASE PARAGRAPHS REBUILT 2026-08-07 against a COMPLETE per-participant breakdown.
% ★ STANDING RULE (author, 2026-08-07): NEVER draft a results claim from one participant
%   unless the passage is explicitly about that participant (as the S01 paragraph is). Run
%   every subject, every time, before writing anything about results. This paragraph was
%   first drafted from S04 alone and understated the cohort badly as a result.
%
% COMPLETE TABLE — each participant's own most-boring vs own most-engaging episode.
% "+" = the measure moves as the participant's report predicts.
%   subj   most-boring   most-engaging   gaze     closure    sleep   fatigue-after
%   R2-01  BOR(8)        INT(6)          +2.58    +0.101     +5      +2
%   R2-02  BOR(9)        INT(6)          +2.64    +0.223     +3      +2
%   R2-03  BOR(3)        CLC(8)          +3.75    +0.053     +2      +1
%   R2-04  BOR(7)        INT(7)          +1.53    +0.007     +2       0 tied
%   S01    BOR(6)        INT(7)          +2.65    +0.020     +1      +1
%   S02    BOR(9)        INT(6)          +1.27    +0.127     +6      +3
%   S03    BOR(9)        CLC(7)          +0.20    +0.105     +6      +5
%   S04    CLC(9)        BOR(7)          +3.29    -0.033 X   +5      +8
%   S05    BOR(9)        CLC(5)          -3.40 X  +0.662     +8      +3
%   S06    INT(8)        BOR(5)          -7.75 X  -0.056 X    0 tied -2 X
%   S07    BOR(9)        CLC(7)          +4.65    +0.360     +6      +4
%   S08    BOR(9)        INT(4)          +3.19    +0.471     +3      +2
% TALLY: gaze 10 agree / 2 contra. closure 10 / 2. sleep 11 / 1 tied / 0 contra.
%        fatigue-after 10 / 1 tied / 1 contra.
% S05's gaze figure is computed from the 22% of their episode where their eyes were open (closure
%   0.7782), which is why it is unreliable rather than contrary.
% S04's closure miss is 3.3 percentage points with BOTH episodes under 9% -- inside the range
%   where both count as attentive watching.
% S06 is the only participant contradicted by more than one measure.
% ---------------------------------------------------------------------------

Taken one participant at a time, the measures follow what each person said about their own
episodes with few exceptions. Setting each participant's self-rated most boring episode against
their self-rated most engaging one, gaze deviation runs in the predicted direction for ten of
the twelve, eye closure for ten, the sleep item for eleven with one tie and no contradiction,
and fatigue afterward for ten with one tie. No measure is carried by a handful of participants
while failing on the rest.

The three exceptions are worth examining rather than counting, because two of them have
explanations that do not depend on the result and one does not. The first is S04, whose ratings
ran opposite to the rest of the cohort: S04 rated the tutorial the most engaging of their three
episodes and the clinical footage the most boring, where ten of the other eleven rated the
tutorial their most boring. Sorted by what S04 reported, three of the four measures follow.
S04's gaze wandered 3.29 degrees further during the episode S04 called boring, S04 fought sleep
at six against one, and S04 finished the episode called boring at nine out of nine for fatigue
against one out of nine for the episode called engaging, the widest fatigue margin recorded in
the study. Only eye closure runs the other way, and the size of that difference is what
matters: 4.8% of the clinical episode against 8.2% of the tutorial, a gap of 3.3 percentage
points. Both figures sit near the bottom of the cohort, whose median is 9.5%, so S04 kept their
eyes open through both films and the measure had almost nothing to separate. One measure with
little room to move runs against S04; three with room to move run with S04.

The second is S05, who spent 78% of the boring film with their eyes shut. S05's closure margin
is by far the largest in the study, and S05's sleep and fatigue scores follow S05's reports.
Only gaze runs contrary, and gaze is the measure S05's own data disqualifies: where the eyes are
closed there is no gaze to record, so S05's figure is computed from the remaining 22% of the
episode. A number built from a fifth of an episode and a number built from a whole one are not
the same measurement, and the disagreement is between those two things rather than between S05
and their eyes.

The third exception is S06, the true anomaly of the study: the one participant whose data does
not track the subjective reports on any measure. S06 rated the film selected to interest
as the most boring of the three, at eight out of nine, and rated the tutorial the most engaging.
Set S06's own most-boring episode against S06's own most-engaging one and every measurement
either contradicts S06 or says nothing. S06's gaze was 7.75 degrees tighter during the episode
S06 called boring, where the cohort's gaze widens. S06's eyes were closed less during it. S06
reported being less tired afterward, not more, and gave the same sleep score for both. Nothing
here is offered as an explanation, because the data does not supply one; S06 is reported as an
exception and left standing as one.

S06 has a consequence for every other number in this section, and it is worth stating openly
rather than leaving for a reader to discover. Every association reported here reaches
significance with S06 included---that is the finding, and it stands on the full cohort of
twelve. What changes when S06 is set aside is not whether the results hold but how emphatically.
The correlation between gaze deviation and reported boredom rises from $+$0.595 to $+$0.706, and
between fatigue afterward and reported boredom from $+$0.640 to $+$0.767; every other
association rises comparably and every probability falls.\footnote{The full comparison, with and
without S06, is set out in the appendix: all eight correlations, the four within-participant
tallies, and what each looks like at twelve participants and at eleven. It makes an interesting
demonstration of how much a single case can hold down a result at this sample size.} S06 is
retained regardless. Nothing depends on the removal, since every association already clears the
threshold at twelve, and the only thing that would have marked S06 for exclusion is that S06's
data disagrees---which is a preference rather than a criterion. The figures reported in this
section are therefore the conservative figures, and the relationships described are stronger
than they appear.

What these cases confirm is the choice of organizing variable, and it is worth stating plainly
rather than leaving implied. Sorted by which film was playing, three of the twelve participants
contradict the cohort, and a film-by-film analysis would record them as noise diminishing an
effect, since their episodes would be averaged in against reports that ran the other way. Sorted by what each participant reported, two of those three cease to be
exceptions at all and take their place among the confirming cases, while the third becomes
something more useful than noise: a single identifiable participant whose measurements do not
follow their own account, available to be examined rather than averaged into the result. The
organizing variable is not a presentational preference. It determines whether the study's most
interesting participants are evidence or interference.

% ---------------------------------------------------------------------------
% BATCH 3 = the S01 case (author-approved 2026-08-07 after verification).
% Figures verified this session against out/features-pooled-wide.csv, out/survey-episodes.csv,
% and a per-sample pass over the Round 1 crop files using the layout-aware reader:
%   S01 BOR self-report: boredom 6/9, engagement 3/9, sleep-fight 2, minutes-until-bored 5.
%   S01 BOR eye closure 0.0895 -- 17th percentile of the cohort's self-rated MOST-BORING
%     episodes, 67th percentile of their MOST-ENGAGING ones.
%   Sleep-fight 2 is the LOWEST of the nine participants reporting boredom >= 6 on that film
%     (others 6,7,7,8,8,9,9,9). Rank 1 of 9; probability by chance 1/9 = 0.111 -- NOT significant,
%     and the prose must not imply otherwise.
%   Minutes-until-bored 5 is the longest of that group (others 0.42, 0.5, 0.5, 2, 2, 3, 4, 4).
%   ★ PER-SAMPLE GAZE CONCENTRATION on the boring film, all 8 R1 participants verified
%     (each cell's closure fraction reproduces the analysis exactly):
%       S08 5.615 deg | S01 8.788 | S05 9.510 | S02 10.619 | S07 10.855 | S03 11.302 |
%       S04 11.741 | S06 14.181   (median offset from that participant's own centre)
%     S05 and S08 are EXCLUDED from the comparison: their eyes were shut for 78% and 58% of
%     the episode, so concentration computed on the fragments left over is an artifact.
%     Among the six with eyes substantially open, S01 is the tightest.
%   S01 held gaze within 10 deg of their own centre 55.8% of the time; S04 40.4%.
% ★ NOT CLAIMED: that S01 and S04 looked at the SAME screen locations. Their eye-in-head
%   directions sit 18.2 deg apart with 24% distributional overlap, but head pose was never
%   instrumented, so eye-in-head cannot be converted to screen position. The author's own
%   observation of the overlay recordings is the basis for any same-location claim, and is
%   attributed as such. Plan for the video analysis that WOULD settle it, should it ever be
%   wanted: PLAN-S01-S04-GAZE-VIDEO-2026-08-07.md.
% ---------------------------------------------------------------------------

S01 complicates the account in the opposite direction, and the complication is
instructive. S01 reported the tutorial as genuinely boring---six of nine, with engagement at
three---and then behaved throughout as though it were not. S01's eyes were shut for 9% of the
episode, which places that episode among the cohort's engaged episodes rather than its bored episodes. S01 reported the least fighting of sleep of any participant who called that film boring,
scoring two where the other eight in that group scored between six and nine, and S01 took five
minutes to become bored, the longest of that group, where several others reported boredom
setting in within a minute.

What separates S01 from everyone else is the gaze. Measured frame by frame across the
episode, S01's eye stayed closer to its own habitual direction than that of any other
participant who watched the film with the eyes meaningfully open, holding within ten degrees of
that direction for 56% of the episode.\footnote{S05 and S08 show tighter figures still and are
set aside rather than counted, having spent 78% and 58% of the episode with the eyes shut. The
concentration measure is then computed on whatever fragments remain open, and brief openings
cluster tightly by construction, so concentration is interpretable only for a participant who
actually watched. Among the six who did, the ordering runs from S01's 8.8 degrees to the widest
at 14.2.} That is tighter than S04, who rated the same film the most engaging of the three and
who held within ten degrees for 40% of it. S04, enjoying the tutorial, looked around it more
than S01, who was bored by it.

Read against the definition the study was built on, this is anomalous. Boredom there is a
failure to engage attention, so a bored viewer should show attention going somewhere else. This
participant's attention went nowhere else. What the case suggests instead is that the capacity to
hold attention is separable from whether the material rewards it---that a person can find
something dull and continue to attend to it closely, and that the two facts can be measured
independently and can disagree. Neither the self-report nor the eye measures can be called wrong
here. They are reporting different things, and the divergence is the finding.

The case will not carry statistical weight and is not offered as though it could. Being the
lowest of nine on the sleep item establishes nothing on its own. What the case does is mark a
possibility the cohort-level results cannot address, and identify the participant in whom it
could be examined further: namely, the power of attention span, and the possibility, if not the
likelihood, that eye tracking can be used to identify and measure it.

% ---------------------------------------------------------------------------
% BATCH: the order control (outline §L3.10 / D-24) and the EEG statement (§L3.8).
% Drafted 2026-08-07 while the D-41 composite re-run was in progress. Neither depends on
% the composite, so neither is affected by that change.
%
% ORDER CONTROL, verified against the pre-D-41 outputs (these figures do not move with the
% composite change since they are raw per-cell values):
%   S01 saw BORING second and INTERESTING last: closure BOR 0.0895 > INT 0.0698.
%   S05 saw the same order:                     closure BOR 0.7782 > INT 0.1394.
%   In both, the LAST-VIEWED film carries the LOWER closure, which is the opposite of what a
%   fatigue-accumulation account predicts.
%   Cohort: boring sits above interesting on closure in 12 of 12.
%   ★ AT n = 2 THIS IS A CHECK, NOT AN ARGUMENT, AND THE PROSE SAYS SO (D-24).
%
% EEG: 0 of 5 surfaces reach significance; no pairwise comparison survives correction; the
%   best surface is default-mode power at Friedman p = 0.1969 with 6 of 8 moving the same way
%   on the two contrasts involving the boring film. R1-only, N=8, never repeated.
%   Framing must match L1: a direction worth noting, not a result. Never per-channel
%   significance for a channel that does not reach it (register v8 standing rule).
% ---------------------------------------------------------------------------

Ten of the twelve participants watched the boring film last, which raises an obvious question
about whether the eye measures are tracking boredom or simply tiredness accumulating across an
afternoon. Two participants happen to answer it. They were the first two through the protocol,
before the viewing order was fixed, and they saw the boring film second and the interesting
film last. If the pattern were an artifact of sequence, their last-viewed film should show the
most eye closure. It does not. In both cases the film they watched last carries the lower
figure, and the boring film watched earlier carries the higher one---0.09 against 0.07 for the
first participant, and 0.78 against 0.14 for the second, a difference of a kind no ordering
effect would produce. Across the whole cohort the boring film sits above the interesting one on
eye closure in twelve cases out of twelve, so the pattern does not depend on the two who saw a
different sequence. At two participants this is a check rather than an argument, and it is
offered as one: it cannot rule out an order effect, only demonstrate that the effect is not
manufactured by the order in the only two cases where the two can be told apart.

The EEG produced nothing this section can use, and the shape of the failure is worth recording
because it is not a failure of equipment. Five measures were derived from the recordings, and
none distinguishes what participants reported at conventional significance. The channel is also the study's thinnest: it was recorded in the
first round only, at eight participants, and was not carried into the second.

% ---------------------------------------------------------------------------
% THE DIVERGENCE DECOMPOSITION (the rebuilt keystone). Drafted 2026-08-07 on the D-41
% composite: eye closure + gaze deviation (per-participant baseline). Pupil is OUT.
%
% Verified against out/keystone-decomposition.csv after the D-41 run:
%   BOR  9 concordant-bored, 3 reverse                       (0 divergent, 0 concordant-engaged)
%   CLC  5 divergent, 4 concordant-engaged, 3 concordant-bored
%   INT  4 divergent, 7 concordant-engaged, 1 concordant-bored
%   ALL 36: divergent 9 | concordant-engaged 11 | concordant-bored 13 | reverse 3
%   ★ Every one of the 12 tutorial composites is NEGATIVE -- no participant's eye measures read
%     as engaged on that film, including the three who did not report boredom (S04, S06, R2-03,
%     who are the 3 'reverse' cells).
% Sensitivity sweep on the clinical film (out/keystone-composite-sensitivity.csv):
%   eye closure alone  -> 7 divergent
%   gaze alone         -> 3 divergent
%   both (baseline)    -> 5 divergent
%   Both measures independently produce divergence in the same direction; no weighting choice
%   creates it.
% ★ FILM-NEUTRAL FRAMING per the governing ruling: the count across all 36 episodes leads;
%   the concentration on clinical and interest films is an observation, not the frame.
% ---------------------------------------------------------------------------

Everything so far has treated the two kinds of evidence separately. Setting them against each
other episode by episode gives the section its central result. For each of the thirty-six
episodes the two eye measures are combined into a single figure, positive where the eyes behaved
as they did during that participant's more engaged episodes and negative where they behaved as
during the less engaged episodes, and that figure is placed beside what the participant said. Four
outcomes are possible: both say engaged, both say bored, or they disagree in one of two
directions.

Twenty-four of the thirty-six episodes agree. Thirteen have both the eyes and the participant
reporting boredom; eleven have both reporting engagement. In nine episodes the eyes read as
engaged while the participant reported boredom. In the remaining three the eyes read as
disengaged while the participant reported no boredom.

The nine matter most, because they are the shape this section set out to look for.
A participant whose eyes hold steady and stay open through an episode they afterward describe as
boring is not producing a contradiction that one measure has to win. The eyes report what the
body did; the participant reports how the episode was to live through; and the two can differ
without either being wrong.

Whether the divergence depends on how the two measures were combined can be tested directly, and
it does not. Taking eye closure by itself produces the divergence in seven cases on the clinical
footage; taking gaze deviation by itself produces it in three; the two combined produce five.
Each measure yields the pattern on its own, and the combination sits between them rather than
manufacturing something neither contains. No weighting decision creates this result, which is
the objection any composite invites.

% ★ P27 CUT 2026-08-09 (author approved). The deleted paragraph claimed that all twelve
% tutorial composites being negative was a finding. VERIFICATION SHOWED IT IS NEARLY AN
% ARTIFACT: the composite is a within-subject z, so each participant's three values sum to
% EXACTLY zero (verified: +0.000000 for all twelve), which makes a participant's worst episode
% negative BY CONSTRUCTION. The tutorial was the worst episode for 11 of 12 participants, so
% "all twelve negative" follows almost automatically. It is also a film-ranking claim, which
% the governing ruling excludes.
%   What survived verification: R2-03 alone is a clean body-versus-report case on that film --
%   closure z -1.40 AND gaze z -1.40, both measures agreeing the episode was the worst of the
%   three, against a reported boredom of 3 and engagement of 6. S06 adds nothing beyond the
%   anomaly already discussed. S04 is NOT a body-versus-report case at all: closure z -1.41
%   against gaze z +0.31, so the two measures disagree with EACH OTHER and the composite is
%   negative only because closure outweighs gaze.
%   If R2-03 is ever wanted, it belongs with the case paragraphs, not in a paragraph about a film.
%
% ★★ FIGURES — AUTHOR INSTRUCTION 2026-08-09: THIS IS THE HEART OF THE RESULTS AND NEEDS
%    MULTIPLE FIGURES PLACED THROUGH IT, not one summary chart at the end. Insert figures in
%    and around these paragraphs so a reader can see the results rather than only read them.
%    Candidates, in the order the prose introduces them:
%      1. Spearman/Pearson comparison chart for every measure against both criteria
%         (already flagged as PENDING earlier in this module).
%      2. Per-participant scatter of gaze deviation against reported boredom, 36 episodes,
%         within-subject z, so the individual points and the fitted relationship are visible.
%      3. The same for eye closure.
%      4. The within-participant paired plots -- own most-boring vs own most-engaging -- for
%         gaze, closure, sleep-fight and fatigue-after, showing the 10/10/11/10 tallies as
%         connected pairs rather than as counts.
%      5. The four-category decomposition across all 36 episodes.
%      6. The with/without-S06 comparison (or hold this one for the appendix).
%    matplotlib is now installed and the 72 tree figures regenerate, so these can be produced.
